\documentclass[10pt,letterpaper]{article}
\usepackage[T1]{fontenc}
\usepackage{lmodern,amsmath,amssymb,amsthm,microtype}
\usepackage[margin=0.9in]{geometry}
\usepackage[hidelinks]{hyperref}
\newtheorem{theorem}{Theorem}
\title{Five colors for a field and its rotated copy}
\author{Aafi Khan Malek\\\small Independent Researcher, Ahmedabad, India\\
\small \href{mailto:aafimalek2032@gmail.com}{aafimalek2032@gmail.com}}
\date{November 2026}
\begin{document}
\maketitle
\begin{abstract}
Let \(K=\mathbb Q(\sqrt{33},\sqrt{-3})\) and
\(\rho=(7+\sqrt{-15})/8\), embedded in \(\mathbb C\) with the
positive imaginary square roots. We construct a proper five-coloring
of the entire induced unit-distance graph on \(K\cup\rho K\).
For its \(2\)-adically integral part a fifth color suffices to separate
one valuation layer. At negative valuations, each low-side integral
coset sees a single color on its higher-valued cross neighbors,
so it can reserve that color and use the other four internally.
\end{abstract}

\section*{The place at two and the cross equation}
Choose a square root of \(33\) in \(\mathbb Q_2\) (an odd
integer congruent to \(1\) modulo \(8\) is a square in
\(\mathbb Q_2\)). Embed \(K\) into \(\mathbb Q_2(\zeta)\) with
\(\zeta^2+\zeta+1=0\), sending \(\sqrt{-3}\) to \(1+2\zeta\).
The residue field is \(\mathbb F_4\); let \(v\) be the additive
valuation with \(v(2)=1\) and \(O=\{x\in K:v(x)\geq0\}\).
Indeed, in the basis \(1,\zeta\) a common power of \(2\)
can be removed and multiplication of the nonzero normalized
residues remains nonzero in \(\mathbb F_4\).
The automorphism \(\zeta\mapsto-1-\zeta\) preserves valuation,
so complex conjugation on \(K\) does as well. Thus
\(v(N(x))=2v(x)\) for \(N(x)=x\overline x\) and \(x\ne0\).
The residue map \(O/2O\cong\mathbb F_4\) is onto since \(\zeta\)
has all four residues in its \(\mathbb F_2\)-span with \(1\).

Write \(j=\sqrt{-15}\). It is not in \(K\): the square classes
of \(-3,33,-15\) are independent, by the primes \(11,5\)
and then the sign of \(-3\). In particular the two copies of
\(K\) share just \(0\). For \(x,y\in K\), expand
\[
 |x-\rho y|^2=N(x)+N(y)
 -\tfrac78(x\overline y+y\overline x)
 +\tfrac j8(x\overline y-y\overline x).
\]
A rational squared distance forces \(x\overline y=y\overline x\),
since \(1,j\) are independent over \(K\). For nonzero
representatives this gives a real ratio \(t=y/x\) in
\(\mathbb Q(\sqrt{33})\), and the expression reduces to
\[
 |x-\rho y|^2=N(x-y)+\tfrac14 x\overline y
             =N(x)\left(1-\tfrac74t+t^2\right).
 \tag{1}
\]
The first reduced expression is symmetric in \(x,y\). In valuation
arguments write \(a=v(x)\), \(b=v(y)\), \(c=v(x-y)\), allowing \(c=\infty\).
The two terms have valuations \(2c\) and \(a+b-2\).

\section*{The two integral copies}
Initially give each point of \(O\) its residue color, with
the same rule on \(\rho O\) after inverse rotation. Within a side
a unit displacement has valuation zero, hence distinct
endpoint residues. A cross unit edge with equal residues has
\(c\ge1\); by (1) its product term must have valuation zero,
so \(a+b=2\). If \(a\ne b\), then \(c=\min(a,b)=0\),
a contradiction. Thus \(a=b=1\).

More generally, every cross unit incident to a representative
of valuation one has both representatives of valuation one.
When \(a=1,b=0\), the product valuation \(-1\) dominates the
difference valuation \(0\). When \(a=1,b\ge2\), both term
valuations are positive. Neither gives squared length one.
An endpoint of valuation one cannot be unit-separated from \(0\).
The valuation-one set of either side is independent, since its
differences have valuation at least one. Recolor that entire set
on one side with a fifth color. This resolves all equal-residue
cross unit edges and creates no monochromatic edges.

The integral coloring also admits an optional equality for any
specified pair at Euclidean distance two. For a pair in one side,
recolor the opposite side: its same-side difference has valuation one,
so its residues agree. A genuine cross pair of squared length four
cannot have \(a=b=0\) or \(a=b=1\), by (1): the product term
would dominate with valuation \(-2\) or \(0\), respectively.
Choose a marked endpoint of valuation at least one so that its
partner is not of valuation one. This is possible by choosing the
unique valuation-one endpoint if present. Recolor the shell on
the chosen side, then give both terminals the fifth color. The
same-side terminal has no edge to that shell by valuation, and
the opposite-side terminal has none by the cross incidence
classification just proved. The terminals are nonadjacent.

\section*{All negative valuations}
\begin{theorem}
The induced unit-distance graph on \(K\cup\rho K\) is five-colorable.
\end{theorem}
\begin{proof}
Use the preceding integral coloring as the initial condition.
Consider a cross unit edge with minimum valuation \(m<0\).
There is no such edge from \(0\) to a point of negative valuation.
If both nonzero representatives had valuation \(m\), the product
term in (1) would have valuation \(2m-2\) strictly below that
of the other term, which cannot yield a unit. Otherwise interchange
the roles of the representatives if needed, writing the lower
one as \(x\). Then \(c=a=m<b\). As the first term valuation
is \(2m<0\), the second must match it, so \(b=m+2\).

In \(\mathbb Q_2\) the quadratic \(1-\tfrac74t+t^2\)
has roots \(r_-,r_+\), because \(-15\equiv1\bmod 8\).
Their valuations are \(-2,2\): the numerators \(7\pm\sqrt{-15}\)
sum to \(14\) of valuation \(1\) and multiply to \(64\) of
valuation \(6\). Now \(t=y/x\) has valuation two, and
\(v(t-r_-)=-2\). Equation (1) gives
\[
  v(t-r_+)=2-2m,\qquad v(y-r_+x)=2-m\ge3.
 \tag{2}
\]
Here these valuations are taken in the completion,
whether or not \(r_+\) belongs to \(K\).

For a fixed low-side coset \(C=x_0+O\) of valuation \(m\),
every higher cross neighbor, viewed in the opposite native copy,
has valuation \(m+2\). All such neighbors are in a single class
modulo \(2O\): if \(x,x'\in C\), then
\(v(r_+(x-x'))\ge2\), and each associated neighbor has error
of valuation at least three in (2). At integral levels the
base coloring is constant on that class at its fixed valuation
(\(0\) or \(1\)); at negative levels we will give each such
class a single color.

Color successively the two copies' cosets at levels \(m=-1,-2,\ldots\).
For each such coset, its higher neighbors have already been colored
and, if there are any, have one common color. Use four distinct
colors other than that color for its four residue classes modulo
\(2O\); choose any four if there are no higher neighbors.
Each native unit edge lies within one \(O\)-coset and joins distinct
residues. Any cross unit edge with a negative valuation is proper
by the palette reserved at its lower end. Together with the base
case these cover all edges.
\end{proof}

The rule can be made effective without determining whether a coset
has a neighbor. Choose \(x_0\) by truncating the two \(2\)-adic
coordinates to their finite negative binary digits. Approximate
\(r_+\) by a rational \(r'\) to precision \(2-m\);
then \(y_*=r'x_0\) has valuation \(m+2\) and every actual higher
neighbor differs from \(y_*\) by an element of \(2O\).
Reserve the already defined color of \(y_*\), whether or not it is
a neighbor. This recursion strictly increases valuation until it
reaches the integral coloring.

\paragraph{Scope.}
The upper bound applies only to the stated union, not its additive
closure or the entire plane. For this union the bound is sharp by
the known five-chromatic type-M construction: Parts's 509-vertex
example uses two native subgraphs in \(K\) with the same rotation
\(\rho\) (see \href{https://arxiv.org/html/2010.12665v2}{Parts,
arXiv:2010.12665v2, Sections 2, 4 and 6}). We use that published
construction for the lower bound, not a new obstruction.
The companion program checks the rule on an exact 105-point fixture,
including nine rational constructions of cross unit edges with negative
valuation. The finite check is illustrative, not the infinite proof.
The distance-two equality result above is for the integral union
only. This written proof was developed with AI assistance and has
not been formalized in Lean.
\end{document}
